% ==============================================================================
% capitulos/07_consideracoes_finais.tex
% Capítulo 7 — Considerações Finais e Trabalhos Futuros
% ==============================================================================

\chapter{Considerações Finais e Trabalhos Futuros}
\label{cap:consideracoes}

\section{Síntese dos Resultados de Projeto}
\label{sec:cons_sintese}

Este documento técnico apresentou a arquitetura, a especificação e a implementação de referência de uma pipeline híbrida de análise de dados industriais integrando o controlador lógico programável \textbf{WEG PLC500}, o protocolo \textbf{Modbus TCP/IP} e a plataforma \textbf{Node-RED}, desenvolvida como entrega de uma \textbf{Situação de Aprendizagem (SA) de IoT de escopo livre}. A definição desse escopo representou uma evolução pragmática e metodologicamente validada a partir da concepção inicial do projeto --- que previa um servidor de tempo NTPv4 baseado em antena GNSS e microcontrolador ESP32, mas esbarrou na inviabilidade física de recepção satelital em laboratório \textit{indoor}. A realização da pivotagem mediante \textbf{aprovação prévia do professor} garantiu a integral conformidade do projeto com o estabelecido na atividade acadêmica, resultando em uma bancada de Internet Industrial das Coisas (IIoT) autocontida, reproduzível e independente de infraestrutura externa de radiofrequência.

Utilizando quatro botoeiras (simulando sensores de presença, nível crítico, sobretemperatura e parada de emergência NF) e quatro sinalizadores LED (simulando motor, válvula de alívio, alarme e indicador de operação normal), a arquitetura projetada demonstra como conciliar ferramentas de fluxo de eventos e modelos de IA com requisitos estritos de engenharia de controle:
\begin{itemize}
  \item \textbf{Determinismo preservado na camada OT:} a lógica em \textit{Structured Text} (IEC 61131-3) no CODESYS V3.5 garante filtragem de \textit{debounce}, contagem precisa de eventos, supervisão de enlace via \textit{watchdog} de \textit{heartbeat} e intertravamento local \textit{fail-safe} com precedência absoluta sobre comandos remotos.
  \item \textbf{Interoperabilidade transparente via Modbus TCP/IP e Node-RED:} o uso do pacote \texttt{node-red-contrib-modbus} aliado a um mapa de registradores base zero com palavra de retorno das saídas físicas (\texttt{wMascara\_DO\_Real}) assegura auditabilidade completa da malha fechada.
  \item \textbf{Inteligência analítica governada na camada IT:} a combinação entre estatística determinística em janelas deslizantes ($\lambda_i$, $D_i$ e \textit{Z-score}) em nós \texttt{function} e diagnóstico contextual por agente LLM protegido por \textit{JSON Schema} e política de \textit{fallback} evita o acoplamento direto de respostas probabilísticas a atuadores físicos, além de automatizar relatórios técnicos e a rastreabilidade em PostgreSQL.
  \item \textbf{Rigor metodológico na validação:} o protocolo formal de dez casos de teste, os critérios quantitativos de aceitação e a matriz de riscos oferecem um roteiro objetivo para o comissionamento da bancada.
\end{itemize}

\section{Limitações Conhecidas e Trabalhos Futuros}
\label{sec:cons_futuros}

Por empregar exclusivamente entradas e saídas digitais simuladas por botoeiras e LEDs nesta primeira etapa, o projeto concentra-se na dinâmica de eventos discretos e na validação arquitetural da malha OT/IT. Como desdobramentos naturais para trabalhos futuros, destacam-se:

\begin{enumerate}
  \item \textbf{Execução da campanha experimental em bancada:} implantar o sistema conforme o Apêndice~\ref{ap:checklist}, executar os casos \texttt{CT-01} a \texttt{CT-10} e registrar as medições reais de bancada (distribuição de latência RTT nos nós do Node-RED, tempo de resposta do intertravamento no CODESYS Trace e avaliação qualitativa dos diagnósticos emitidos pelo LLM).
  \item \textbf{Expansão para grandezas analógicas reais:} incorporar cartões de expansão da série MOD (compatíveis com a linha RUW100 do WEG PLC500) para leitura de sinais contínuos ($4\text{--}20\,\text{mA}$, $0\text{--}10\,\text{V}$ e termorresistências PT100), ampliando a camada analítica no Node-RED com técnicas multivariadas (\textit{Isolation Forest} ou EWMA).
  \item \textbf{Fortalecimento de cibersegurança no transporte:} avaliar a migração da comunicação supervisória para \textbf{OPC~UA com assinatura digital e criptografia TLS} --- suportado tanto no CODESYS do WEG PLC500 quanto no ecossistema Node-RED --- reduzindo a superfície de exposição associada ao Modbus TCP/IP em redes ampliadas \cite{stouffer_nist_ot, iec_62443}.
  \item \textbf{Observabilidade visual e sincronismo de tempo na rede OT:} integrar painéis gráficos diretamente à base PostgreSQL (via Grafana ou Node-RED Dashboard) e, caso futuramente se disponha de acesso a antena externa em topo de edifício, reintegrar o conceito original do servidor NTP Stratum~1 (ESP32 + GNSS) como fonte local de sincronismo de tempo para o WEG PLC500 e para o servidor Node-RED na sub-rede isolada \cite{rfc5905_ntp}.
\end{enumerate}
